\documentclass[11pt,a4paper]{article}

\usepackage[margin=2.5cm]{geometry}
\usepackage{amsmath,amssymb,amsthm}
\usepackage{booktabs}
\usepackage{graphicx}
\usepackage[colorlinks=true,linkcolor=blue,citecolor=blue]{hyperref}

\newtheorem{lemma}{Lemma}
\newtheorem{remark}{Remark}
\newtheorem{sanity}{Consistency check}
\newtheorem{warning}{Validity warning}

\newcommand{\RR}{\mathbb{R}}
\newcommand{\EE}{\mathbb{E}}
\newcommand{\PP}{\mathbb{P}}
\newcommand{\Var}{\mathrm{Var}}
\newcommand{\Cov}{\mathrm{Cov}}
\newcommand{\sd}{\mathrm{sd}}
\newcommand{\Phibar}{\overline{\Phi}}
\newcommand{\htwo}{h_2}
\newcommand{\ksub}{k_{\mathrm{sub}}}
\newcommand{\Ebase}{E_{\mathrm{base}}}
\newcommand{\Emax}{E_{\mathrm{max}}}
\newcommand{\uhat}{\hat{u}}

\title{Mutual information of a homomer receptor array \\ in a single-family chemical environment}
\author{Analytical baseline for the octopus chemosensation project}
\date{\today}

\begin{document}
\maketitle

\begin{abstract}
\noindent
We compute, in closed form up to one-dimensional quadratures, the mutual information
(MI) between the binary activity pattern of an array of $R$ freely placed receptors
(homomers) and the chemical stimulus, for a single Gaussian family of ligands in a
$D$-dimensional chemical latent space, one ligand per sniff, and a shared lognormal
concentration. The derivation uses the \emph{gaussian affinity kernel} of the simulator
throughout, and introduces the quadratic kernel only as an explicitly flagged limit. The
answer is governed by a single dimensionless number $\nu$, the ratio of the common-mode
threshold fluctuation to the private per-receptor signal. The array delivers its ideal
$R$ bits when $\nu \ll 1$, degrades as $R/\nu$ when $\nu \gtrsim 1$, and collapses to
$O(\log R)$ bits when $\nu \gtrsim R/\log R$. Because the per-unit baseline energy
$\Ebase$ and selectivity ceiling $\Emax$ are both learnable, $\nu$ is not an environment
parameter but an outcome the optimizer partly controls, and the saturation of the
gaussian kernel is what stops it from being driven to zero. Heteromers are deferred
entirely to Appendix~\ref{app:hetero}.
\end{abstract}

\section{Overview: approach, key steps, and logic}
\label{sec:overview}

\paragraph{The problem this solves.}
In the numerical work so far, the environment was tuned so that a homomer array reaches
its intrinsic architectural optimum, $R$ bits for $R$ binary receptors. That was the
right choice for isolating architecture from geometry, but it has a cost: once the
homomer arm is pinned at the ceiling, every parameter sweep is flat and there is nothing
to compare against. To study the concentration channel with a single ligand per sniff we
need to leave the ceiling in a controlled way, and that requires a theory predicting what
a homomer array actually delivers in a given environment, not an upper bound on what it
could deliver. This document is that theory.

\paragraph{The approach in one paragraph.}
The calculation rests on one structural observation. A receptor fires when the
log-concentration exceeds its half-activation energy, and with the gaussian kernel that
condition rearranges exactly into \emph{the match between ligand and receptor exceeds a
demand level set by the concentration}. The match is $\exp(-d^2/\lambda^2)$ with $d$ the
latent distance, so taking a logarithm turns the comparison into an inequality on $d^2$,
and expanding the square splits the ligand dependence into two pieces of completely
different character: a term $\|l\|^2$ that is \emph{the same for every receptor}, and a
term $u_r\cdot l$ that is \emph{linear in the ligand coordinate and specific to each
receptor}. The concentration joins the first group, because one ligand at one
concentration presents the same concentration to every receptor. The whole array is
therefore driven by a single scalar common mode plus $R$ Gaussian projections, and
projections of an isotropic Gaussian onto orthogonal directions are independent, so the
array reduces to \emph{$R$ independent Gaussian variables cut by one common random
threshold}. Conditionally on that threshold the receptors are independent and
identically distributed Bernoulli variables, the joint response is exchangeable, and its
entropy follows from a two-line identity. Everything then depends on the environment
through a single ratio.

\paragraph{The logic, step by step.}
\begin{enumerate}
  \item \textbf{Match versus demand.} Rewrite the activation condition exactly, with the
        gaussian kernel, as a comparison between a geometric match and a concentration
        demand. Show the firing set is exactly a ball whose radius is set by the
        concentration (Section~\ref{sec:step1}).
  \item \textbf{Separation.} Expand the square and split the condition into a shared
        scalar and a private linear projection (Section~\ref{sec:step2}).
  \item \textbf{Symmetry.} Isotropy forces equal receptor norms at a balanced optimum,
        and orthogonal directions make the projections independent
        (Section~\ref{sec:step3}).
  \item \textbf{Conditional independence.} Condition on the shared threshold. The
        receptors become independent coins with a common bias $q$
        (Section~\ref{sec:step4}).
  \item \textbf{Statistics of the common mode, kernel-free.} Mean, variance, balance
        condition, and the definition of the control parameter $\nu$, all written for an
        arbitrary kernel (Section~\ref{sec:step5}).
  \item \textbf{Specialisation to the gaussian kernel.} Evaluate the common mode
        explicitly, identify the three learnable knobs and what each controls, and
        recover the quadratic kernel as the small-field limit
        (Section~\ref{sec:gaussian}).
  \item \textbf{The optimal lever arm.} The genuinely new content that the gaussian
        kernel brings: moving receptors outward suppresses the common mode linearly but
        amplifies concentration noise exponentially, so there is a finite optimum
        (Section~\ref{sec:lever}).
  \item \textbf{Entropy.} Convert conditional independence into an exact MI expression,
        by two routes that must agree (Section~\ref{sec:step6}).
  \item \textbf{Limits.} The three regimes, with explicit constants and a worked example
        (Section~\ref{sec:limits}).
  \item \textbf{Error control.} Quantify the one approximation made, and give the exact
        reduction that avoids it (Section~\ref{sec:exact}).
  \item \textbf{Decomposition.} Split the MI into identity and concentration parts
        analytically (Section~\ref{sec:split}).
  \item \textbf{Symmetry breaking.} Release the equal-norm assumption and show a
        staggered array beats a balanced one precisely when $\nu > 1$
        (Section~\ref{sec:stagger}).
  \item \textbf{The other corner.} What changes when $R > D$
        (Section~\ref{sec:corner}).
\end{enumerate}

\paragraph{What the reader should take away before any algebra.}
Four statements carry the document. First, a shared log-concentration fluctuation and a
radial breathing of the ligand cloud enter at exactly the same place and are
information-theoretically the same object. Second, for an array of balanced receptors
that common mode is pure nuisance, so concentration spread \emph{reduces} the total MI
even though it increases input entropy. Third, because the baseline energy $\Ebase$ is
learnable, the array has a free parameter whose only function is to defeat the common
mode, and the saturation of the gaussian kernel is the sole reason it cannot do so
completely: this is why the quadratic kernel is degenerate here and why the gaussian one
must be used. Fourth, converting the common mode into information requires staggering
the receptor thresholds, which is free once $\Ebase$ is learnable.

\section{Model, notation, and assumptions}
\label{sec:model}

\subsection{Notation}

Every symbol used in the main text is defined here. Vectors are column vectors in
$\RR^D$ and $\|\cdot\|$ is the Euclidean norm. Entropies and mutual informations are in
bits.

\begin{center}
\begin{tabular}{lll}
\toprule
Symbol & Meaning & Type \\
\midrule
$D$ & dimension of the chemical latent space & integer \\
$R$ & number of receptors in the array & integer \\
$l \in \RR^D$ & latent coordinate of the ligand presented in a sniff & random \\
$\sigma$ & per-coordinate standard deviation of the ligand cloud & parameter \\
$c$, $s = \ln c$ & concentration and log-concentration & random \\
$m$, $\sigma_c$ & mean and standard deviation of $s$ & parameters \\
$\lambda$ & affinity length scale of the gaussian kernel & parameter (fixed) \\
$u_r \in \RR^D$ & latent coordinate of receptor $r$ & learnable \\
$U_r = \|u_r\|$, $\uhat_r = u_r/U_r$ & norm and direction of receptor $r$ & learnable \\
$\Ebase^{(r)}$ & baseline energy, i.e.\ $\ln EC_{50}$ at a perfect match & learnable \\
$\Emax^{(r)}$ & selectivity ceiling, extra energy at zero match & learnable, $>0$ \\
$d_r = \|l - u_r\|$ & latent distance from ligand to receptor $r$ & random \\
$K_r = e^{-d_r^2/\lambda^2} \in (0,1]$ & \textbf{match} of the ligand to receptor $r$ & random \\
$Z_r$ & \textbf{demand}, Eq.~\eqref{eq:demand} & random \\
$\bar{z}_r$ & mean demand, the operating point of receptor $r$ & derived \\
$\Gamma_r = \Emax^{(r)}/\sigma_c$ & \textbf{dynamic range} of receptor $r$ & derived \\
$a_r \in \{0,1\}$, $A$, $K$ & activity, array pattern, number active & random \\
$p_r = \uhat_r\cdot l$ & private projection of the ligand on receptor $r$ & random \\
$\Theta$ & common mode, Eq.~\eqref{eq:Theta} & random \\
$\tau$ & shared threshold on the private projections & random \\
$q$ & firing probability conditional on $\tau$ & random \\
$\nu$ & control parameter, Eq.~\eqref{eq:nudef} & derived \\
$\rho = \sigma\sqrt{D}/\lambda$ & cloud width over kernel range & derived \\
$w = U^2/\lambda^2$ & squared lever arm in kernel units & derived \\
$\phi$, $\Phi$, $\Phibar$ & standard normal density, lower tail, upper tail & functions \\
$\htwo(p)$ & binary entropy in bits & function \\
$W(\cdot)$ & Lambert $W$ function, the inverse of $x \mapsto xe^x$ & function \\
\bottomrule
\end{tabular}
\end{center}

\subsection{The model}

\paragraph{Environment.} One chemical family, an isotropic Gaussian cloud centred at the
origin of the latent space (putting the centre at the origin is a choice of coordinates
and costs nothing, since only differences $l - u_r$ appear):
\begin{equation}
  l \sim \mathcal{N}(0,\ \sigma^2 I_D).
  \label{eq:ligandlaw}
\end{equation}
Each sniff presents exactly one ligand from \eqref{eq:ligandlaw} at a log-concentration
drawn independently,
\begin{equation}
  s = \ln c \sim \mathcal{N}(m,\ \sigma_c^2),
  \qquad s \perp l .
  \label{eq:conclaw}
\end{equation}

\paragraph{Receptors and the gaussian kernel.} Receptor $r$ is specified by three
learnable objects: a position $u_r \in \RR^D$, a baseline energy $\Ebase^{(r)} \in \RR$,
and a selectivity ceiling $\Emax^{(r)} > 0$. Its half-activation concentration follows
the \textbf{gaussian affinity kernel}, which is the
\texttt{affinity\_kernel="gaussian"} branch of the simulator,
\begin{equation}
  \boxed{\;\ln EC_{50}^{(r,l)}
   = \Ebase^{(r)} + \Emax^{(r)}\left(1 - e^{-\|l-u_r\|^2/\lambda^2}\right). }
  \label{eq:kernel}
\end{equation}
It is called saturating because as the distance grows the energy tends to the finite
ceiling $\Ebase + \Emax$ instead of diverging. The alternative
\texttt{affinity\_kernel="quadratic"} branch, $\ln EC_{50} = \Ebase + dE\,\|l-u_r\|^2$,
is recovered in Section~\ref{sec:gaussian} as the small-field limit of \eqref{eq:kernel}
and is \emph{not} assumed anywhere before that.

\paragraph{Readout.} The receptor is a deterministic threshold device,
\begin{equation}
  a_r = \mathbf{1}\!\left[\, s > \ln EC_{50}^{(r,l)} \,\right].
  \label{eq:readout}
\end{equation}

\subsection{Assumptions, listed so they can be audited}

\begin{description}
  \item[A1 (gaussian kernel, exact).] Equation~\eqref{eq:kernel} is used as written. No
    expansion of the kernel is made in Sections~\ref{sec:step1} to~\ref{sec:step5}. The
    quadratic kernel appears only in Section~\ref{sec:gaussian}, clearly labelled, and
    Section~\ref{sec:lever} shows that it is \emph{degenerate} for the question of the
    optimal receptor placement, which is the reason the gaussian form is carried
    throughout.
  \item[A2 (one ligand per sniff).] The stimulus is a single pair $(l,s)$. This is the
    regime where chemical geometry rather than mixture combinatorics limits the array.
    Mixtures are a separate calculation.
  \item[A3 (deterministic readout).] The map $(l,s)\mapsto A$ is deterministic and
    noiseless, so $H(A\mid l,s)=0$ and
    \begin{equation}
      I\big(A;(l,s)\big) = H(A) - H(A\mid l,s) = H(A).
      \label{eq:MIisH}
    \end{equation}
    Every ``MI'' below means $H(A)$. With a soft readout at finite temperature $T$ an
    extra term $\EE[\sum_r \htwo(a_r)]$ must be subtracted.
  \item[A4 (independence of the common mode).] In Sections~\ref{sec:step4}
    to~\ref{sec:limits} the shared threshold is treated as independent of the private
    projections. This is the only genuine approximation in the main derivation. Its error
    is $O(1/D)$ per receptor and it is quantified, and removed, in
    Section~\ref{sec:exact}.
  \item[A5 (orthogonal directions).] Sections~\ref{sec:step3} onward assume $R\le D$ and
    mutually orthogonal receptor directions. The opposite corner is
    Section~\ref{sec:corner}.
  \item[A6 (identical, balanced receptors).] Sections~\ref{sec:step3}
    to~\ref{sec:split} assume all receptors share the same $U$, $\Ebase$ and $\Emax$, and
    sit at the balanced operating point. This is justified by isotropy in
    Section~\ref{sec:step3} and deliberately released in Section~\ref{sec:stagger}.
  \item[A7 (delta method for the demand).] Section~\ref{sec:gaussian} expands
    $\ln Z$ around its mean. Its validity condition is stated explicitly there, and
    Section~\ref{sec:lever} shows the optimum sits at the edge of that condition, so the
    resulting optimum is correct in scaling but not in constant.
\end{description}

\section{Step 1: match versus demand, and the exact ball}
\label{sec:step1}

\paragraph{Step 1a. Name the geometric factor.} The only place the ligand coordinate
enters the kernel \eqref{eq:kernel} is through
\begin{equation}
  K_r \;\equiv\; \exp\!\left(-\frac{\|l-u_r\|^2}{\lambda^2}\right) \;\in\; (0,1],
\end{equation}
which we call the \textbf{match} between the ligand and receptor $r$. It equals $1$ for
a perfectly docked ligand and tends to $0$ for a chemically remote one. In terms of it,
\eqref{eq:kernel} reads
\begin{equation}
  \ln EC_{50}^{(r,l)} = \Ebase^{(r)} + \Emax^{(r)}\,(1 - K_r),
\end{equation}
so the half-activation energy slides linearly between $\Ebase$ at perfect match and
$\Ebase + \Emax$ at zero match.

\paragraph{Step 1b. Rearrange the firing condition.} Insert this into the readout
\eqref{eq:readout} and isolate $K_r$. Each step is elementary:
\begin{align}
  a_r = 1
  &\iff s > \Ebase + \Emax(1-K_r)
   &&\text{(definition)} \notag\\
  &\iff s - \Ebase - \Emax > -\Emax K_r
   &&\text{(expand and move terms)} \notag\\
  &\iff \Emax K_r > \Ebase + \Emax - s
   &&\text{(multiply by }-1) \notag\\
  &\iff K_r > \frac{\Ebase + \Emax - s}{\Emax}
   &&(\Emax > 0).
\end{align}
Define the \textbf{demand}
\begin{equation}
  \boxed{\;Z_r \;\equiv\; \frac{\Ebase^{(r)} + \Emax^{(r)} - s}{\Emax^{(r)}}
         \;=\; 1 - \frac{s - \Ebase^{(r)}}{\Emax^{(r)}},
  \qquad\text{so that}\qquad
  a_r = \mathbf{1}\!\left[\,K_r > Z_r\,\right]. }
  \label{eq:demand}
\end{equation}
\textbf{A receptor fires when the match exceeds the demand.} This is exact, no expansion
has been made, and it is the cleanest statement of what the gaussian kernel does.

\paragraph{Step 1c. The demand is Gaussian, and dimensionless.} Since $Z_r$ is an affine
function of $s$ and $s$ is Gaussian by \eqref{eq:conclaw},
\begin{equation}
  Z_r \sim \mathcal{N}\!\left(\bar z_r,\ \frac{1}{\Gamma_r^{2}}\right),
  \qquad
  \bar z_r = 1 - \frac{m - \Ebase^{(r)}}{\Emax^{(r)}},
  \qquad
  \Gamma_r = \frac{\Emax^{(r)}}{\sigma_c}.
  \label{eq:Zlaw}
\end{equation}
Two learnable quantities have acquired a clean meaning, and they do different jobs:
\begin{itemize}
  \item $\Ebase$ sets $\bar z$, the \textbf{operating point}: how demanding the receptor
    is on average, hence its firing rate. This is the balance knob.
  \item $\Emax$ sets $\Gamma$, the \textbf{dynamic range}: the span of
    log-concentrations the receptor can resolve, measured in units of the concentration
    spread. Large $\Gamma$ means a receptor whose demand barely moves when the
    concentration fluctuates, that is, a receptor insensitive to concentration noise.
\end{itemize}

\paragraph{Step 1d. The three regimes, and where saturation lives.} Because the match is
confined to $(0,1]$, the demand being Gaussian is unbounded and two saturated outcomes
occur with nonzero probability:
\begin{center}
\begin{tabular}{lll}
\toprule
Demand & Outcome & Probability \\
\midrule
$Z_r \le 0$ & fires whatever the ligand, carries nothing & $\Phi(-\bar z_r \Gamma_r)$ \\
$0 < Z_r \le 1$ & fires if and only if $l$ lies in a ball, Step~1e & $1$ minus the others \\
$Z_r > 1$ & never fires, carries nothing & $\Phibar\big((1-\bar z_r)\Gamma_r\big)$ \\
\bottomrule
\end{tabular}
\end{center}
These two atoms are the entire content of the word ``saturating''. They are absent from
the quadratic kernel, where the demand is unbounded in both directions, and
Section~\ref{sec:lever} shows that their absence is exactly what makes the quadratic
kernel degenerate.

\paragraph{Step 1e. The informative regime is an exact ball.} For $0 < Z_r \le 1$ take
the logarithm of \eqref{eq:demand}. The logarithm is increasing so the inequality
direction is preserved, and $\ln K_r = -\|l-u_r\|^2/\lambda^2$:
\begin{equation}
  K_r > Z_r
  \iff -\frac{\|l-u_r\|^2}{\lambda^2} > \ln Z_r
  \iff \boxed{\;\|l - u_r\|^2 \;<\; \mathrm{Rad}_r^2 \;\equiv\; -\lambda^2 \ln Z_r \;}
  \label{eq:ball}
\end{equation}
and $-\ln Z_r \ge 0$ on this range, so the radius is real. \textbf{The receptive field of
a homomer is exactly a ball in latent space, for the gaussian kernel as much as for the
quadratic one.} The only thing the kernel decides is how the radius depends on the
concentration:
\begin{equation}
  \text{gaussian:}\quad \mathrm{Rad}^2(s) = \lambda^2 \ln\!\frac{\Emax}{\Ebase+\Emax-s},
  \qquad\qquad
  \text{quadratic:}\quad \mathrm{Rad}^2(s) = \frac{s - \Ebase}{dE}.
  \label{eq:radtwo}
\end{equation}
This is worth pausing on, because it is what keeps the calculation cheap. Everything
from Step~2 to Step~9 treats $\mathrm{Rad}^2(s)$ as an unspecified scalar function of a
scalar, so the whole derivation is kernel-agnostic and the kernel is substituted once,
at the end.

\section{Step 2: separating the common mode from the private signal}
\label{sec:step2}

\paragraph{Step 2a. Expand the square.} This is the only algebra that matters, so it is
written in full. For any two vectors,
\begin{align}
  \|l-u_r\|^2
   &= (l-u_r)\cdot(l-u_r) \notag\\
   &= l\cdot l - 2u_r\cdot l + u_r\cdot u_r \notag\\
   &= \underbrace{\|l\|^2}_{\text{same for all } r}
    \;-\; \underbrace{2u_r\cdot l}_{\text{linear in } l,\ \text{specific to } r}
    \;+\; \underbrace{\|u_r\|^2}_{\text{constant, specific to } r}.
  \label{eq:expand}
\end{align}

\paragraph{Step 2b. Insert into the ball condition.} Using \eqref{eq:ball},
\begin{equation}
  a_r = 1
  \iff \|l\|^2 - 2u_r\cdot l + U_r^2 < \mathrm{Rad}_r^2
  \iff 2u_r\cdot l > \|l\|^2 + U_r^2 - \mathrm{Rad}_r^2 .
\end{equation}

\paragraph{Step 2c. Factor the norm out of the linear term.} Write $u_r = U_r\uhat_r$
with $\|\uhat_r\|=1$, so $u_r \cdot l = U_r p_r$ with
\begin{equation}
  p_r \;\equiv\; \uhat_r \cdot l
\end{equation}
the \textbf{private projection}: the coordinate of the ligand along receptor $r$'s own
direction. Dividing by the positive number $2U_r$, which preserves the inequality,
\begin{equation}
  \boxed{\;a_r = \mathbf{1}\!\left[\,p_r > \tau_r\,\right],
  \qquad
  \tau_r = \frac{\Theta_r + U_r^2}{2U_r},
  \qquad
  \Theta_r \equiv \|l\|^2 - \mathrm{Rad}_r^2(s) . }
  \label{eq:Theta}
\end{equation}

\paragraph{Step 2d. Why $\Theta$ is called the common mode.} Under A6 all receptors share
$\Ebase$ and $\Emax$, so $\mathrm{Rad}_r^2 = \mathrm{Rad}^2$ and $\Theta_r = \Theta$
carries no index. It is one scalar random variable that the whole array is compared
against, and it contains two contributions that enter identically:
\begin{itemize}
  \item $\|l\|^2$, how far the presented ligand sits from its family centre, a purely
    geometric radial coordinate;
  \item $-\mathrm{Rad}^2(s)$, minus the receptive-field radius, which is a monotone
    decreasing function of the log-concentration.
\end{itemize}

\begin{remark}[A shared concentration and a breathing cloud are the same object]
This is the first headline statement, and it is visible already in \eqref{eq:Theta}. A
fluctuation of $\ln c$ shared by all receptors, and a fluctuation of the radius $\|l\|$
of the ligand cloud, appear in the same scalar with the same sign structure, so they
have identical information-theoretic consequences for the array. Everything proved below
about the concentration spread is automatically a statement about the radial spread of
the chemical family, and conversely.
\end{remark}

\begin{remark}[The two distinct roles of the receptor norm]
In \eqref{eq:Theta} the norm $U_r$ appears twice, and the two effects differ:
\begin{equation}
  \frac{\partial \tau_r}{\partial \Theta} = \frac{1}{2U_r},
  \qquad
  \tau_r\big|_{\Theta=0} = \frac{U_r}{2}.
\end{equation}
A distant receptor is \emph{less} sensitive to the common mode, which is good, and has a
\emph{higher} baseline threshold, which lowers its firing rate. Section~\ref{sec:lever}
is about this trade-off, and Section~\ref{sec:stagger} uses it again.
\end{remark}

\section{Step 3: what isotropy forces on the design}
\label{sec:step3}

\begin{lemma}[The marginal firing rate depends on a receptor only through its norm]
\label{lem:norm}
Under \eqref{eq:ligandlaw} and \eqref{eq:conclaw}, and at fixed $\Ebase$ and $\Emax$,
the probability $\PP(a_r=1)$ is a function of $U_r=\|u_r\|$ alone and is strictly
decreasing in it.
\end{lemma}

\begin{proof}
By \eqref{eq:ball}, $a_r = \mathbf{1}[\|l-u_r\|^2 < \mathrm{Rad}^2(s)]$ with
$\mathrm{Rad}^2$ a function of $s$ only, and $s \perp l$. Since
$l \sim \mathcal{N}(0,\sigma^2 I_D)$, the vector $(l-u_r)/\sigma$ is normal with mean
$-u_r/\sigma$, so
\begin{equation}
  \frac{\|l-u_r\|^2}{\sigma^2}\sim\chi'^2_D\!\left(\frac{\|u_r\|^2}{\sigma^2}\right),
\end{equation}
a noncentral chi-square with $D$ degrees of freedom and noncentrality
$\|u_r\|^2/\sigma^2$. Its law depends on $u_r$ only through $\|u_r\|$, giving the first
claim. A noncentral chi-square is stochastically increasing in its noncentrality, so
$\|l-u_r\|^2$ stochastically increases with $U_r$ and the probability that it falls below
an independent threshold strictly decreases.
\end{proof}

\begin{remark}[Why this forces equal norms, given equal $\Ebase$ and $\Emax$]
A binary variable has maximal entropy at marginal $1/2$. By Lemma~\ref{lem:norm}, at
fixed $\Ebase$ and $\Emax$ the equation $\PP(a_r=1)=1/2$ has exactly one solution in
$U_r$, the same for every receptor, so all receptors sit on one sphere and only their
directions are free. That is assumption A6. Note the caveat that matters later: once
$\Ebase$ is allowed to differ between receptors, balance can be achieved at \emph{any}
norm by compensating with the baseline, and the argument for equal norms collapses. That
is the subject of Sections~\ref{sec:lever} and \ref{sec:stagger}.
\end{remark}

\begin{lemma}[Orthogonal directions give independent private projections]
\label{lem:orth}
If $\uhat_1,\dots,\uhat_R$ are orthonormal and $R\le D$ then $p_1,\dots,p_R$ are
independent with $p_r\sim\mathcal{N}(0,\sigma^2)$. In general
$\Cov(p_r,p_{r'})=\sigma^2\,\uhat_r\cdot\uhat_{r'}$.
\end{lemma}

\begin{proof}
$p=\hat V^\top l$ with $\hat V$ the $D\times R$ matrix of directions. A linear image of a
Gaussian is Gaussian, with mean $0$ and covariance
$\hat V^\top(\sigma^2 I_D)\hat V = \sigma^2\hat V^\top \hat V$, whose entries are
$\sigma^2\uhat_r\cdot\uhat_{r'}$. Orthonormality makes this $\sigma^2 I_R$, and for a
Gaussian vector a diagonal covariance is equivalent to independence. Extending
$\{\uhat_r\}$ to an orthonormal basis exhibits the $p_r$ as coordinates of $l$ in that
basis, which are independent and identically distributed $\mathcal{N}(0,\sigma^2)$.
\end{proof}

Under A5 and A6 the array therefore collapses to its final and much simpler form,
\begin{equation}
  \boxed{\;a_r = \mathbf{1}[p_r > \tau],
  \qquad
  \tau = \frac{\Theta + U^2}{2U},
  \qquad
  p_r \stackrel{\text{iid}}{\sim}\mathcal{N}(0,\sigma^2),}
  \label{eq:tau}
\end{equation}
with one and the same $\tau$ for every receptor: \emph{$R$ independent Gaussian variables
cut by a single common random threshold}.

\section{Step 4: conditioning on the common mode}
\label{sec:step4}

\paragraph{Step 4a. Freeze the common mode.} Adopt A4 and treat $\tau$ as independent of
$(p_1,\dots,p_R)$. Conditionally on $\tau=t$ the activities $a_r=\mathbf{1}[p_r>t]$ are
functions of independent variables, hence independent, and identically distributed:
\begin{equation}
  q(t) \;\equiv\; \PP(p_r > t)
       \;=\; \PP\!\left(\frac{p_r}{\sigma} > \frac{t}{\sigma}\right)
       \;=\; \Phibar\!\left(\frac{t}{\sigma}\right).
  \label{eq:q}
\end{equation}

\paragraph{Step 4b. The conditional law of the pattern.} For a pattern
$a\in\{0,1\}^R$ with $k=\sum_r a_r$ active receptors,
\begin{align}
  \PP(A=a\mid \tau=t)
   &= \prod_{r=1}^R q(t)^{a_r}\big(1-q(t)\big)^{1-a_r}
     &&\text{(conditional independence)}\notag\\
   &= q(t)^{\sum_r a_r}\big(1-q(t)\big)^{R-\sum_r a_r}
     &&\text{(collect equal factors)}\notag\\
   &= q(t)^{k}\big(1-q(t)\big)^{R-k}.
     &&\text{(definition of } k)
\end{align}
The right-hand side depends on the pattern \emph{only through $k$}, which is the
property that makes the problem tractable.

\paragraph{Step 4c. Average over the common mode.} Writing $\Pi_k$ for the resulting
pattern probability,
\begin{equation}
  \PP(A=a) = \EE_\tau\!\left[q(\tau)^k\big(1-q(\tau)\big)^{R-k}\right] \equiv \Pi_k,
  \qquad k=\textstyle\sum_r a_r .
  \label{eq:Pi}
\end{equation}
So $A$ is an \textbf{exchangeable} binary vector, a mixture of independent and
identically distributed Bernoulli vectors with mixing variable $q(\tau)\in[0,1]$. This is
the finite de Finetti representation, obtained constructively rather than invoked.

\paragraph{Step 4d. The law of the count.} There are $\binom{R}{k}$ patterns with $k$
active receptors, each of probability $\Pi_k$, so
\begin{equation}
  \PP(K=k) = \binom{R}{k}\Pi_k, \qquad k=0,\dots,R.
  \label{eq:PK}
\end{equation}

\begin{sanity}
$\sum_k \binom{R}{k}\Pi_k
 = \EE_\tau\big[\sum_k\binom{R}{k}q^k(1-q)^{R-k}\big]
 = \EE_\tau\big[(q+1-q)^R\big] = 1$ by the binomial theorem. This is the first assertion
to put in any implementation.
\end{sanity}

\section{Step 5: statistics of the common mode, for any kernel}
\label{sec:step5}

This section uses only that $\mathrm{Rad}^2$ is some function of $s$, so it holds for
both kernels.

\paragraph{Step 5a. Mean and variance of $\|l\|^2$.} Write
$\|l\|^2=\sum_{i=1}^D l_i^2$ with $l_i$ independent $\mathcal{N}(0,\sigma^2)$. Then
$l_i^2/\sigma^2\sim\chi^2_1$, which has mean $1$ and variance $2$. Summing $D$
independent copies and multiplying by $\sigma^2$,
\begin{equation}
  \EE\big[\|l\|^2\big]=D\sigma^2, \qquad \Var\big[\|l\|^2\big]=2D\sigma^4 .
  \label{eq:radial}
\end{equation}

\paragraph{Step 5b. Mean and variance of the common mode.} From
$\Theta=\|l\|^2-\mathrm{Rad}^2(s)$ with $s\perp l$, variances add:
\begin{equation}
  \EE[\Theta]=D\sigma^2-\EE[\mathrm{Rad}^2],
  \qquad
  \Var[\Theta]=2D\sigma^4+\Var[\mathrm{Rad}^2].
  \label{eq:Thetastats}
\end{equation}
Equation~\eqref{eq:Thetastats} is the quantitative version of the remark in
Section~\ref{sec:step2}: the geometric term and the concentration term sit side by side
in one sum.

\paragraph{Step 5c. The balance condition.} Apply the affine map
$\tau=(\Theta+U^2)/(2U)$ and set $\EE[\tau]=0$, which centres the firing probability on
$q=1/2$. The numerator must vanish:
\begin{equation}
  \boxed{\;\EE\big[\mathrm{Rad}^2(s)\big] \;=\; D\sigma^2 + U^2 . }
  \label{eq:balance}
\end{equation}
In words, and independently of the kernel: \emph{the mean squared radius of the
receptive field must equal the mean squared distance from the receptor to a ligand},
since $\EE\|l-u_r\|^2 = D\sigma^2 + U^2$. Note what has changed relative to a model with
a fixed baseline. Equation~\eqref{eq:balance} constrains $\Ebase$, which is learnable,
and therefore it does \emph{not} determine $U$. The norm remains free, and fixing it is
the job of Section~\ref{sec:lever}.

\begin{remark}[Exact symmetry of the balanced point]
At $\EE[\tau]=0$ with $\tau$ symmetric about zero the mean firing rate is exactly one
half: $\tau \overset{d}{=} -\tau$ gives
$\EE[\Phibar(\tau/\sigma)]=\EE[\Phibar(-\tau/\sigma)]=\EE[1-\Phibar(\tau/\sigma)]$,
hence $2\EE[q]=1$. Balance controls the mean of $q$, not merely its median.
\end{remark}

\paragraph{Step 5d. The control parameter.} All that remains is the \emph{spread} of
$q$, and by \eqref{eq:q} that spread is the standard deviation of $\tau$ measured in
units of the private-projection scale $\sigma$:
\begin{equation}
  \boxed{\;\nu \;\equiv\; \frac{\sd(\tau)}{\sigma}
   \;=\; \frac{\sqrt{\,2D\sigma^4+\Var[\mathrm{Rad}^2(s)]\,}}{2U\sigma}, }
  \qquad\text{so that}\qquad
  \frac{\tau}{\sigma} = \nu z,\quad z\sim\mathcal{N}(0,1),\quad q=\Phibar(\nu z),
  \label{eq:nudef}
\end{equation}
the last form holding at balance and to the extent that $\Theta$ is Gaussian, which it is
for large $D$ by the central limit theorem applied to $\|l\|^2$. In words, $\nu$ is the
ratio of the common-mode fluctuation, which pushes every receptor the same way, to the
private signal, which pushes each receptor differently.

\section{Step 6: specialising to the gaussian kernel}
\label{sec:gaussian}

\paragraph{Step 6a. The radius in terms of the demand.} From \eqref{eq:ball},
$\mathrm{Rad}^2=-\lambda^2\ln Z$ with $Z\sim\mathcal{N}(\bar z,\Gamma^{-2})$ by
\eqref{eq:Zlaw}. So the two moments we need are those of $\ln Z$.

\paragraph{Step 6b. Expand the logarithm about the mean demand (A7).} Write
$Z=\bar z(1+\varepsilon)$ with $\varepsilon=(Z-\bar z)/\bar z$, so $\varepsilon$ has mean
zero and standard deviation $1/(\Gamma\bar z)$. Using
$\ln(1+\varepsilon)=\varepsilon-\varepsilon^2/2+O(\varepsilon^3)$,
\begin{equation}
  \EE[\ln Z] \simeq \ln\bar z - \frac{1}{2\Gamma^2\bar z^2},
  \qquad
  \Var[\ln Z] \simeq \frac{1}{\Gamma^2\bar z^2}.
  \label{eq:deltamethod}
\end{equation}

\begin{warning}
\label{warn:delta}
Equation~\eqref{eq:deltamethod} requires $\Gamma\bar z \gg 1$, that is, a demand whose
fluctuations are small compared to its mean. Equivalently it requires the saturation
probabilities of Step~1d to be negligible. This is not automatic, because the optimizer
will be shown to drive $\bar z$ down. Step~7e evaluates the condition \emph{at the
optimum} in closed form and finds that it holds comfortably for $\rho\lesssim1$ and
fails for $\rho\gtrsim1.5$, so the formulas below are quantitative in the first regime
and only indicative in the second.
\end{warning}

\paragraph{Step 6c. Impose balance and read off the operating point.} Substituting
$\EE[\mathrm{Rad}^2]=-\lambda^2\EE[\ln Z]\simeq-\lambda^2\ln\bar z$ into the balance
condition \eqref{eq:balance}, and introducing the two natural dimensionless groups
\begin{equation}
  \rho^2 \equiv \frac{D\sigma^2}{\lambda^2}
   \quad(\text{cloud width over kernel range, squared}),
  \qquad
  w \equiv \frac{U^2}{\lambda^2}
   \quad(\text{squared lever arm in kernel units}),
\end{equation}
we obtain
\begin{equation}
  -\lambda^2\ln\bar z = D\sigma^2+U^2
  \qquad\Longleftrightarrow\qquad
  \boxed{\;\bar z = e^{-(\rho^2+w)} . }
  \label{eq:zbar}
\end{equation}
The mean demand must equal the typical match, which is what balance means in the
language of Step~1b. Since $\bar z$ is set by $\Ebase$ through \eqref{eq:Zlaw}, this
determines the learnable baseline and nothing else.

\paragraph{Step 6d. The concentration contribution to the common mode.} Combining
\eqref{eq:deltamethod} and \eqref{eq:zbar},
\begin{equation}
  \sd\big[\mathrm{Rad}^2\big] = \lambda^2\,\sd[\ln Z]
   \simeq \frac{\lambda^2}{\Gamma\bar z}
   = \frac{\lambda^2}{\Gamma}\,e^{\rho^2+w}.
  \label{eq:sdrad}
\end{equation}
\textbf{The concentration noise seen by the receptor grows exponentially with its
squared lever arm.} The mechanism is transparent in the match-versus-demand picture: a
distant receptor needs a large receptive field to stay balanced, a large field means a
small demand $\bar z$, and a small demand is a small number whose logarithm is violently
sensitive to the Gaussian jitter of size $1/\Gamma$.

\paragraph{Step 6e. The control parameter for the gaussian kernel.} Substitute
\eqref{eq:sdrad} into \eqref{eq:nudef}. Using $\sigma^2/\lambda^2=\rho^2/D$ to make
everything dimensionless,
\begin{equation}
  \boxed{\;\nu^2(w) \;=\; \frac{1}{w}\left[\,\underbrace{\frac{\rho^2}{2}}_{\textstyle P}
   \;+\; \underbrace{\frac{D\,e^{2\rho^2}}{4\rho^2\Gamma^2}}_{\textstyle Q}\,e^{2w}\right]
   \;=\; \frac{P+Qe^{2w}}{w} . }
  \label{eq:nugauss}
\end{equation}
The first term is the radial common mode, which a longer lever arm suppresses. The
second is the concentration common mode, which a longer lever arm amplifies. Their
competition is the content of the next section.

\paragraph{Step 6f. The quadratic kernel as the small-field limit.} Take
$\rho^2+w\ll 1$, meaning both the ligand cloud and the receptor offset are small
compared with the kernel range, so by \eqref{eq:zbar} $\bar z\to 1$ and the receptive
field radius $\mathrm{Rad}^2 \ll \lambda^2$. Then
\begin{equation}
  \mathrm{Rad}^2 = -\lambda^2\ln Z
   = -\lambda^2\ln\!\left(1-\frac{s-\Ebase}{\Emax}\right)
   \simeq \lambda^2\,\frac{s-\Ebase}{\Emax}
   = \frac{s-\Ebase}{dE},
   \qquad dE \equiv \frac{\Emax}{\lambda^2},
\end{equation}
which is exactly the \texttt{quadratic} branch of \eqref{eq:radtwo}. In the same limit
$e^{2w}\to1$ in \eqref{eq:nugauss} and
\begin{equation}
  \nu^2_{\text{quad}}(w) = \frac{1}{w}\left[\frac{\rho^2}{2}+\frac{D}{4\rho^2\Gamma^2}\right],
  \qquad\text{equivalently}\qquad
  \nu^2_{\text{quad}} = \frac{2D\sigma^4 + \sigma_c^2/dE^2}{4U^2\sigma^2}.
  \label{eq:nuquad}
\end{equation}
The quadratic kernel is therefore not a different model, it is the gaussian kernel for a
receptor whose field is small compared with its own affinity range.

\begin{remark}[Why the quadratic kernel cannot answer the placement question]
Expression \eqref{eq:nuquad} is \emph{monotonically decreasing} in $w$, with no interior
optimum. Under the quadratic kernel the optimizer can send $U\to\infty$, compensate with
the learnable baseline to stay balanced, drive $\nu\to0$, and collect the full $R$ bits
in any environment. That is a degeneracy of the kernel, not a property of the chemistry,
and it is the reason this document carries the gaussian form. It also predicts what one
would see in a simulation run with \texttt{affinity\_kernel="quadratic"}: receptor norms
drifting outward without converging, and baselines drifting down to match.
\end{remark}

\section{Step 7: the optimal lever arm}
\label{sec:lever}

\paragraph{Step 7a. Minimise over the lever arm.} Differentiate \eqref{eq:nugauss} using
the quotient rule:
\begin{equation}
  \frac{\mathrm{d}}{\mathrm{d}w}\,\frac{P+Qe^{2w}}{w}
   = \frac{2Qe^{2w}\,w-\big(P+Qe^{2w}\big)}{w^2},
\end{equation}
so the stationarity condition is
\begin{equation}
  Q\,e^{2w}\,(2w-1) = P .
  \label{eq:stationary}
\end{equation}

\paragraph{Step 7b. Solve it with the Lambert function.} Substitute
$\Lambda \equiv 2w-1$, so $w=(\Lambda+1)/2$ and $e^{2w}=e^{\Lambda+1}$:
\begin{equation}
  Q\,e^{\Lambda+1}\,\Lambda = P
  \quad\Longleftrightarrow\quad
  \Lambda e^{\Lambda} = \frac{P}{Q\,e}
  \quad\Longleftrightarrow\quad
  \boxed{\;\Lambda = W\!\left(\frac{P}{Qe}\right), \qquad w_\star=\frac{1+\Lambda}{2}, }
  \label{eq:lambertsol}
\end{equation}
where $W$ is the Lambert function, defined as the inverse of $x\mapsto xe^x$.

\paragraph{Step 7c. The minimum value.} At the optimum \eqref{eq:stationary} gives
$Qe^{2w_\star}=P/\Lambda$, so
\begin{equation}
  \nu_{\min}^2 = \frac{P+P/\Lambda}{w_\star}
   = \frac{P\,(1+1/\Lambda)}{(1+\Lambda)/2}
   = \frac{2P}{\Lambda}\cdot\frac{(\Lambda+1)/\Lambda}{(\Lambda+1)/\Lambda}
   = \frac{2P}{\Lambda},
\end{equation}
that is, substituting $P=\rho^2/2$,
\begin{equation}
  \boxed{\;\nu_{\min}^2 = \frac{\rho^2}{\Lambda},
   \qquad \Lambda = W\!\left(\frac{2\rho^4\Gamma^2}{D}\,e^{-2\rho^2-1}\right). }
  \label{eq:numin}
\end{equation}

\paragraph{Step 7d. Read it in the large-dynamic-range limit.} For large argument
$W(x)\simeq\ln x-\ln\ln x$, and here
$\ln x = 2\ln\Gamma+\ln(2\rho^4/D)-2\rho^2-1$, so $\Lambda$ grows like $2\ln\Gamma$.
Keeping only that leading term,
\begin{equation}
  \boxed{\;\nu_{\min}\;\simeq\;\frac{\rho}{\sqrt{2\ln\Gamma}},
   \qquad
   w_\star \simeq \ln\Gamma,
   \qquad\text{that is}\qquad
   U_\star \simeq \lambda\sqrt{\ln\Gamma}, }
  \label{eq:scaling}
\end{equation}
with $\rho=\sigma\sqrt{D}/\lambda$ and $\Gamma=\Emax/\sigma_c$. The subleading constants
dropped here are not small at realistic $\Gamma$: for $D=10$, $\rho^2=0.4$ and
$\Gamma=100$ one finds $\Lambda=2.90$ against $2\ln\Gamma=9.21$. Use the Lambert form
\eqref{eq:numin} for numbers and \eqref{eq:scaling} for intuition. Three statements
follow, and they are the practical output of the document.
\begin{enumerate}
  \item \textbf{The common mode can be suppressed, but only logarithmically.} The
    penalty from the ligand cloud is linear in $\rho$, while the array's defence against
    it grows only as $\sqrt{\ln\Gamma}$. Doubling the dynamic range of every receptor
    buys a factor $\sqrt{1+\ln 2/\ln\Gamma}$, which is nothing.
  \item \textbf{The controlling environment parameter is $\rho$, not the concentration
    gap.} The combination $m-\Ebase$ is set by the optimizer through \eqref{eq:zbar}, so
    sweeping it does nothing. The knob is $\rho=\sigma\sqrt{D}/\lambda$, which is the
    same kernel-geometry ratio identified in
    \texttt{doc/theory/08\_environmental\_entropy\_limits.md}.
  \item \textbf{A directly checkable prediction.} Optimised receptors should sit at
    distance $U_\star\simeq\lambda\sqrt{\ln\Gamma}$ from the cloud centre, which is a few
    times $\lambda$ and is \emph{not} the cloud radius $\sigma\sqrt{D}$, together with
    strongly negative baselines given by \eqref{eq:zbar}. This can be read off any
    already-optimised array without running anything new.
\end{enumerate}

\paragraph{Step 7e. Self-consistency: is the delta method valid at its own optimum?}
The validity group of Warning~\ref{warn:delta} can be evaluated at $w_\star$ in closed
form, which is worth doing because the answer is not obvious. Use
$\Lambda e^{\Lambda}=A$ with $A=P/(Qe)$, so $e^{-\Lambda/2}=\sqrt{\Lambda/A}$, and
$\bar z = e^{-(\rho^2+w_\star)}$ with $w_\star=(1+\Lambda)/2$:
\begin{align}
  \Gamma\bar z\Big|_{w_\star}
   &= \Gamma\,e^{-\rho^2}\,e^{-1/2}\,e^{-\Lambda/2}
    = \Gamma\,e^{-\rho^2-1/2}\sqrt{\frac{\Lambda}{A}} \notag\\
   &= \Gamma\,e^{-\rho^2-1/2}\,
      \sqrt{\frac{\Lambda D\,e^{2\rho^2+1}}{2\rho^4\Gamma^2}}
    = \boxed{\;\sqrt{\frac{\Lambda D}{2\rho^4}}\;}
  \label{eq:selfcons}
\end{align}
The dependence on $\Gamma$ cancels exactly: \emph{at its own optimum the receptor's
operating point is set by the geometry alone}. The delta method is therefore
self-consistent whenever $\rho^4 \ll \Lambda D/2$, that is for a ligand cloud narrow
compared with the kernel range, and it degrades as $\rho$ approaches $D^{1/4}$.
Numerically, for $D=10$ and $\sigma=1$:

\begin{center}
\begin{tabular}{rrrrrrr}
\toprule
$\lambda$ & $\rho$ & $\Gamma$ & $\Lambda$ & $U_\star$ & $\nu_{\min}$ & $\Gamma\bar z$ \ (sat.\ prob.) \\
\midrule
$5.0$ & $0.63$ & $10$  & $0.37$ & $4.13$ & $1.04$ & $3.38$ \ ($4\times10^{-4}$) \\
$5.0$ & $0.63$ & $100$ & $2.90$ & $6.98$ & $0.37$ & $9.52$ \ ($\approx 0$) \\
$2.0$ & $1.58$ & $10$  & $0.24$ & $1.58$ & $3.21$ & $0.44$ \ ($0.33$) \\
$1.5$ & $2.11$ & $10$  & $0.02$ & $1.07$ & $15.0$ & $0.07$ \ ($0.47$) \\
\bottomrule
\end{tabular}
\end{center}

Three things to take from the table. First, the first row is an ordinary parameter
choice that lands at $\nu_{\min}\approx1$, the most informative point: the interesting
regime is not exotic. Second, $\lambda$ alone sweeps $\nu$ over more than a decade,
which makes it the knob of choice. Third, saturation is negligible where the theory is
quantitative and becomes dominant only once $\rho\gtrsim1.5$, so the atoms of Step~1d
are not what creates the optimum. What creates it is the \emph{curvature} of
$\mathrm{Rad}^2=-\lambda^2\ln Z$: because the match is bounded above by $1$, a receptor
with a large field must operate at small $\bar z$, where the logarithm amplifies the
concentration jitter. A kernel without that bound, such as the quadratic one, has no
such amplification and no optimum.

\begin{remark}[Saturation is an exact, additive loss]
This can be stated without any expansion. By Step~1d, whenever $Z\le0$ or $Z>1$ the
receptor's conditional firing probability is $1$ or $0$, and $\htwo(1)=\htwo(0)=0$, so
those sniffs contribute nothing to the per-receptor efficiency $\EE[\htwo(q)]$ defined in
Section~\ref{sec:step6}. Hence the exact bound
\begin{equation}
  \EE[\htwo(q)] \;\le\; 1 - \Phi(-\bar z\Gamma) - \Phibar\big((1-\bar z)\Gamma\big),
\end{equation}
valid for any design. The saturation probability is subtracted from the array's
efficiency one for one.
\end{remark}

\section{Step 8: the mutual information, by two routes}
\label{sec:step6}

Two derivations are given because they must agree, and their agreement is the best
available check on an implementation.

\subsection{Route 1: direct summation over patterns}

\paragraph{Step 8a.} Start from the definition and group the $2^R$ patterns by their
number of active receptors, using that $\PP(A=a)=\Pi_k$ depends only on $k$:
\begin{align}
  H(A) &= -\sum_{a\in\{0,1\}^R}\PP(A=a)\log_2\PP(A=a) \notag\\
       &= -\sum_{k=0}^{R}\ \sum_{a:|a|=k}\Pi_k\log_2\Pi_k \notag\\
       &= -\sum_{k=0}^{R}\binom{R}{k}\Pi_k\log_2\Pi_k .
\end{align}

\paragraph{Step 8b.} Substitute $\Pi_k=\PP(K=k)/\binom{R}{k}$ from \eqref{eq:PK} and use
$\log_2(x/y)=\log_2 x-\log_2 y$:
\begin{align}
  H(A) &= -\sum_k \PP(K=k)\log_2\!\left[\frac{\PP(K=k)}{\binom{R}{k}}\right] \notag\\
       &= -\sum_k \PP(K=k)\log_2 \PP(K=k) + \sum_k \PP(K=k)\log_2\binom{R}{k},
\end{align}
that is
\begin{equation}
  \boxed{\;H(A)=H(K)+\EE\!\left[\log_2\binom{R}{K}\right].}
  \label{eq:grouping}
\end{equation}
The first term is the information carried by \emph{how many} receptors fire, the second
by \emph{which} ones, which here is just the logarithm of the number of equally likely
arrangements. This is the same grouping identity already used in the project's
discrete-information notes, so the two formalisms agree.

\subsection{Route 2: the chain rule through the mixing variable}

\paragraph{Step 8c.} Condition on $\tau$ and use $H(A)=H(A\mid\tau)+I(A;\tau)$. By
conditional independence and additivity of entropy over independent variables,
$H(A\mid\tau=t)=\sum_r \htwo(q(t))=R\,\htwo(q(t))$, so
\begin{equation}
  \boxed{\;H(A)=R\,\EE\big[\htwo(q)\big]+I(A;q).}
  \label{eq:chain}
\end{equation}

\paragraph{Step 8d. Reading \eqref{eq:chain}.} The first term is a \emph{per-receptor
yield}: each receptor contributes $\EE[\htwo(q)]\le1$ bits instead of the ideal bit, and
the deficit is caused entirely by the common mode pulling $q$ away from $1/2$. The
second is what the array learns about the common mode itself, which for a one-parameter
mixture grows only logarithmically,
\begin{equation}
  I(A;q)\simeq\tfrac12\log_2 R+\text{const}
  \qquad (R\to\infty,\ \nu \text{ fixed}),
  \label{eq:clarke}
\end{equation}
the standard asymptotic for the mutual information between a scalar parameter and $R$
conditionally independent observations of it. The constant involves the entropy of the
mixing density and the Fisher information of the Bernoulli family. The precise statement
should be verified before it is cited; only the $\frac12\log_2 R$ scaling is used below.

\paragraph{Step 8e. The working formula.} Combining,
\begin{equation}
  \boxed{\;\mathrm{MI}=H(A)\simeq R\,
   \underbrace{\EE_z\big[\htwo\big(\Phibar(\nu z)\big)\big]}_{\textstyle \eta(\nu)}
   +\tfrac12\log_2 R+\text{const},}
  \label{eq:working}
\end{equation}
where $\eta(\nu)\in[0,1]$ is the \textbf{information efficiency per receptor}. It is a
single one-dimensional Gaussian integral, computable to machine precision by
Gauss--Hermite quadrature, and it depends on nothing but $\nu$.

\section{Step 9: limits, asymptotics, and worked numbers}
\label{sec:limits}

\subsection{The ideal limit $\nu\to0$}

\paragraph{Step 9a.} Expand the Gaussian tail about zero. Since $\Phibar'=-\phi$ and
$\Phibar(0)=1/2$,
\begin{equation}
  q=\Phibar(\nu z)=\tfrac12-\phi(0)\nu z+O(\nu^3),
  \qquad \phi(0)=\frac{1}{\sqrt{2\pi}},
\end{equation}
so $\epsilon\equiv q-\tfrac12=-\phi(0)\nu z+O(\nu^3)$.

\paragraph{Step 9b.} Expand the binary entropy about its maximum. Its derivatives are
$\htwo'(p)=\log_2\frac{1-p}{p}$ and $\htwo''(p)=-\frac{1}{\ln2}\frac{1}{p(1-p)}$, so
$\htwo'(1/2)=0$, $\htwo''(1/2)=-4/\ln2$, and
\begin{equation}
  \htwo\!\left(\tfrac12+\epsilon\right)=1-\frac{2}{\ln2}\epsilon^2+O(\epsilon^4).
\end{equation}

\paragraph{Step 9c.} Take the expectation with $\EE[z^2]=1$ and $\phi(0)^2=1/(2\pi)$:
\begin{equation}
  \EE[\epsilon^2]=\frac{\nu^2}{2\pi}
  \qquad\Longrightarrow\qquad
  \boxed{\;\eta(\nu)=1-\frac{\nu^2}{\pi\ln2}+O(\nu^4)=1-0.4592\,\nu^2+O(\nu^4).}
  \label{eq:smallnu}
\end{equation}
A $10\%$ loss per receptor requires $\nu\approx0.47$. Numerically this expansion is good
to about $2\%$ at $\nu=0.5$ and useless beyond $\nu\approx0.7$.

\subsection{The exactly solvable point $\nu=1$}

\paragraph{Step 9d.} At $\nu=1$ the tail function is evaluated at a draw from its own
distribution, and $\Phibar(z)\sim\mathrm{Uniform}(0,1)$. Hence, using
$\htwo(q)=\htwo(1-q)$ to merge the two halves of the binary entropy,
\begin{align}
  \eta(1)&=\int_0^1\htwo(q)\,\mathrm{d}q
         =-2\int_0^1 q\log_2 q\,\mathrm{d}q
         =-\frac{2}{\ln2}\int_0^1 q\ln q\,\mathrm{d}q \notag\\
         &=-\frac{2}{\ln2}\left[\frac{q^2}{2}\ln q-\frac{q^2}{4}\right]_0^1
          =\boxed{\frac{1}{2\ln2}\approx0.7213\ \text{bits}.}
  \label{eq:nu1}
\end{align}
At $\nu=1$ the array has lost $28\%$ of its information per receptor, which is
comfortably measurable. Numerical quadrature confirms $\eta(1)=0.72135$.

\subsection{The common-mode-dominated limit $\nu\gg1$}

\paragraph{Step 9e.} Change variables, $x=\nu z$, so $\mathrm{d}z=\mathrm{d}x/\nu$:
\begin{equation}
  \eta(\nu)=\int\phi(z)\,\htwo\big(\Phibar(\nu z)\big)\mathrm{d}z
           =\frac1\nu\int\phi\!\left(\frac{x}{\nu}\right)\htwo\big(\Phibar(x)\big)\mathrm{d}x .
\end{equation}

\paragraph{Step 9f.} The factor $\htwo(\Phibar(x))$ is appreciable only for
$|x|\lesssim4$, over which $\phi(x/\nu)\simeq\phi(0)$ when $\nu\gg1$. Hence
\begin{equation}
  \boxed{\;\eta(\nu)\simeq\frac{\phi(0)\mathcal{A}}{\nu}\approx\frac{1.0397}{\nu},
  \qquad
  \mathcal{A}=\int_{-\infty}^{\infty}\htwo\big(\Phibar(x)\big)\mathrm{d}x=2.60608 . }
  \label{eq:largenu}
\end{equation}
The constant was obtained by quadrature and the asymptote is accurate to better than
$5\%$ for $\nu\ge3$. The memorable form is $R_{\mathrm{eff}}=R\,\eta(\nu)\approx R/\nu$:
\textbf{$\nu$ counts the receptors the array wastes per receptor it uses.}

\subsection{The three regimes}

\begin{center}
\begin{tabular}{lll}
\toprule
Regime & Mutual information & Character \\
\midrule
$\nu\lesssim1$ & $R(1-0.459\nu^2)$ & near-ideal, flat in every parameter \\
$1\lesssim\nu\lesssim R/\log_2 R$ & $\approx1.04R/\nu+\tfrac12\log_2 R$ & graceful degradation \\
$\nu\gtrsim R/\log_2 R$ & $O(\log_2 R)$, down to $1$ bit & array slaved to the common mode \\
\bottomrule
\end{tabular}
\end{center}
The second crossover follows from $1.04R/\nu=\tfrac12\log_2R$, giving
$\nu_\star\approx2.08R/\log_2R$. The collapse to a single bit is therefore \emph{not} at
$\nu\sim1$: between $\nu=1$ and $\nu\sim R/\log R$ the array degrades as $R/\nu$.

\subsection{Worked numbers for the gaussian kernel}

Combining the optimal lever arm of Step~7e with $\eta(\nu)$, for $D=10$ and $\sigma=1$:
\begin{center}
\begin{tabular}{rrrrrl}
\toprule
$\lambda$ & $\rho$ & $\Gamma=\Emax/\sigma_c$ & $\nu_{\min}$ & $\eta(\nu_{\min})$ & MI per receptor \\
\midrule
$5.0$ & $0.63$ & $100$ & $0.37$ & $0.942$ & near-ideal, sweeps look flat \\
$5.0$ & $0.63$ & $10$  & $1.04$ & $0.706$ & $29\%$ loss, clearly measurable \\
$2.0$ & $1.58$ & $10$  & $3.21$ & $0.308$ & strongly degraded \\
$1.5$ & $2.11$ & $10$  & $15.0$ & $0.069$ & array nearly collapsed \\
\bottomrule
\end{tabular}
\end{center}
The operational message is that \textbf{$\lambda$ is the knob}. Holding the environment
fixed and varying the affinity length scale alone moves the array from the saturated
regime to near-collapse. Sweeping $m-\Ebase$ instead does nothing, because the optimizer
resets the baseline through \eqref{eq:zbar}.

\section{Step 10: what assumption A4 costs, and how to remove it}
\label{sec:exact}

\paragraph{Step 10a. Where the dependence hides.} Decompose the ligand in the orthonormal
receptor basis extended to all of $\RR^D$,
\begin{equation}
  l=\sum_{j=1}^R p_j\uhat_j+l_\perp,
  \qquad
  \|l\|^2=\sum_{j=1}^R p_j^2+W,
  \qquad
  W\equiv\|l_\perp\|^2\sim\sigma^2\chi^2_{D-R}.
\end{equation}
So the common mode $\Theta=\|l\|^2-\mathrm{Rad}^2(s)$ is not independent of the private
projections: it contains their squares. Assumption A4 ignored this.

\paragraph{Step 10b. The exact per-receptor condition.} Since $u_r=U\uhat_r$ has
component $U$ along $\uhat_r$ and zero elsewhere,
$\|l-u_r\|^2=(p_r-U)^2+\sum_{j\neq r}p_j^2+W$, so by \eqref{eq:ball}
\begin{equation}
  \boxed{\;a_r=1 \iff (p_r-U)^2<\mathcal{R}_r^2,
  \qquad
  \mathcal{R}_r^2\equiv \mathrm{Rad}^2(s)-W-\sum_{j\neq r}p_j^2 . }
  \label{eq:exactcond}
\end{equation}
The firing set in a receptor's own private coordinate is an \emph{interval}
$|p_r-U|<\mathcal{R}_r$, not a half-line.

\paragraph{Step 10c. The two things A4 actually did.}
\begin{enumerate}
  \item \textbf{The far edge was dropped.} The condition $p_r<U+\mathcal{R}_r$ was
    discarded. Typically $\mathcal{R}_r\sim\sigma\sqrt{D}$ while
    $p_r\sim\mathcal{N}(0,\sigma^2)$, so the far edge is about $\sqrt{D}$ standard
    deviations away and is crossed with probability $\sim\Phibar(\sqrt{D})$, which is
    $8\times10^{-4}$ at $D=10$ and shrinks fast.
  \item \textbf{The self-term was treated as common.} Putting $p_r^2$ into the shared
    $\Theta$ rather than into receptor $r$'s own radius creates a spurious correlation
    between a receptor and its own threshold. Its weight is the variance share of one
    squared projection,
    \begin{equation}
      \frac{\Var(p_r^2)}{\Var(\Theta)}=\frac{2\sigma^4}{2D\sigma^4+\Var[\mathrm{Rad}^2]}
      \;\le\;\frac1D ,
    \end{equation}
    so the per-receptor error is $O(1/D)$ and the aggregate over the array is $O(R/D)$.
    The cross terms $p_j^2$ with $j\neq r$ are genuinely common and were handled
    correctly.
\end{enumerate}

\paragraph{Step 10d. The exact computation, if it is ever needed.} Condition on the pair
$(W,s)$, which is independent of $p$, and write
$\mathcal{R}^2=\mathrm{Rad}^2(s)-W$. Then \eqref{eq:exactcond} says receptor $r$ fires if
and only if the $R$-dimensional vector $p$ lies in the ball of radius $\mathcal{R}$
centred at $Ue_r$, with $e_r$ the $r$-th basis vector of $\RR^R$. Hence
\begin{equation}
  \PP(A=a)=\EE_{W,s}\left[\gamma_R\!\left(
    \bigcap_{r:a_r=1}B(Ue_r,\mathcal{R})\ \cap
    \bigcap_{r:a_r=0}B(Ue_r,\mathcal{R})^c\right)\right],
\end{equation}
with $\gamma_R$ the $\mathcal{N}(0,\sigma^2I_R)$ measure. The problem in $D+1$ random
dimensions has been reduced exactly to $R+2$, with all the $D$ dependence collapsed into
the scalar $W$. For modest $R$ this is a tractable Monte Carlo and it is the right way to
validate the closed form rather than to replace it.

\section{Step 11: splitting identity from concentration}
\label{sec:split}

The project needs $I(A;\text{identity})$ and $I(A;c)$ separately, and in this model both
are analytic.

\paragraph{Step 11a. Condition on the ligand.} Fix $l$. The only remaining randomness is
$s$, and by \eqref{eq:demand} receptor $r$ fires if and only if $K_r>Z$, where the
matches $K_r$ are now fixed numbers and the demand $Z$ is a monotone decreasing function
of $s$.

\paragraph{Step 11b. The conditional pattern is a thermometer.} Order the receptors by
decreasing match, equivalently by increasing distance from $l$,
\begin{equation}
  K_{(1)}\ge K_{(2)}\ge\dots\ge K_{(R)},
  \qquad
  t_k \;\equiv\; \Ebase+\Emax\big(1-K_{(k)}\big)
\end{equation}
so that $t_1\le t_2\le\dots\le t_R$ are the log-concentrations at which the receptors
switch on, in order. As $s$ increases, receptors turn on nearest-first and in no other
order, so given $l$ the pattern takes only $R+1$ values, ``the $k$ nearest are active'',
with
\begin{equation}
  w_0(l)=F(t_1),
  \qquad
  w_k(l)=F(t_{k+1})-F(t_k),
  \qquad
  w_R(l)=1-F(t_R),
  \qquad
  F(x)=\Phi\!\left(\frac{x-m}{\sigma_c}\right).
  \label{eq:thermo}
\end{equation}

\paragraph{Step 11c. The decomposition.} Therefore
\begin{equation}
  \boxed{\;H(A\mid l)=-\sum_{k=0}^{R}w_k(l)\log_2 w_k(l)\;\le\;\log_2(R+1),
  \qquad
  I(A;l)=H(A)-\EE_l\big[H(A\mid l)\big]. }
  \label{eq:split}
\end{equation}
The identity information is the total minus a term bounded by $\log_2(R+1)$ bits.
Combined with \eqref{eq:working}, in the near-ideal regime
$I(A;l)=R\eta(\nu)-O(\log R)$: almost all of the array's information is about which
ligand was presented, and the concentration channel contributes only logarithmically.
Conditioning on $s$ instead freezes every receptive-field radius simultaneously and
gives $I(A;s)=H(A)-\EE_s[H(A\mid s)]$, with $H(A\mid s)$ the entropy of the cells of an
arrangement of $R$ equal-radius balls under the Gaussian ligand measure.

\begin{remark}[Why this matters for the current numerics]
Equation \eqref{eq:split} replaces, in this limit, the concentration-binning procedure
currently used to estimate $I(A;c)$. It is therefore a target the estimator must
reproduce, which is a far stronger test of the pipeline than internal convergence checks.
\end{remark}

\section{Step 12: releasing A6, and the transition at $\nu=1$}
\label{sec:stagger}

Assumption A6 made every receptor identical. This section releases it. The key fact,
and the reason this section is different from what it would have been with a fixed
baseline, is that $\Ebase$ is learnable per receptor, so \emph{threshold diversity is
free}: it costs no displacement in latent space.

\paragraph{Step 12a. The general conditional-independence formula.} Return to
\eqref{eq:Theta} and keep $U_r$, $\Ebase^{(r)}$, $\Emax^{(r)}$ distinct. The threshold
$\tau_r$ is then a deterministic function of the \emph{two} shared scalars
$\big(\|l\|^2, s\big)$, and conditionally on that pair the receptors are independent
because each depends on its own private projection. Writing
\begin{equation}
  q_r\big(\|l\|^2,s\big)
   = \Phibar\!\left(\frac{1}{\sigma}\,
     \frac{\|l\|^2-\mathrm{Rad}_r^2(s)+U_r^2}{2U_r}\right),
\end{equation}
the chain rule of Step~8c gives the general working formula
\begin{equation}
  \boxed{\;H(A)=\EE_{\|l\|^2,\,s}\!\left[\sum_{r=1}^R \htwo\big(q_r\big)\right]
   +I\big(A;(\|l\|^2,s)\big). }
  \label{eq:general}
\end{equation}
This is a two-dimensional quadrature and it contains everything in the document as
special cases: A6 collapses the pair to the single scalar $\Theta$ and recovers
\eqref{eq:working}.

\paragraph{Step 12b. The two terms pull in opposite directions.} The first term is
maximised by putting every $q_r$ as close to $1/2$ as possible, which means identical
balanced receptors. The second is maximised by making the receptors respond to
\emph{different} values of the shared variables, which means staggered thresholds.
Which dominates is decided by $\nu$.

\paragraph{Step 12c. Identical receptors in the strong-common-mode limit.} With A6 and
$\nu\gg1$ the shared threshold swings over a range far wider than the private
projections, so $q$ is almost always $0$ or $1$ and all receptors fire together or not
at all. From \eqref{eq:grouping}, $K\in\{0,R\}$ with probability $1/2$ each, so $H(K)=1$
and $\log_2\binom{R}{0}=\log_2\binom{R}{R}=0$, giving $H(A)\to1$ bit. An array of $R$
identical receptors has degenerated into a single one.

\paragraph{Step 12d. Staggered receptors in the same limit.} Now give receptor $r$ its
own baseline. For $\nu\gg1$ neglect the private projections, so by \eqref{eq:Theta}
\begin{equation}
  a_r \approx \mathbf{1}\!\left[\,\mathrm{Rad}_r^2(s) > \|l\|^2+U_r^2\,\right].
\end{equation}
The receptors switch on in the order of their thresholds as the shared variables move:
the array has become a \textbf{thermometer code} on the common mode, with $R+1$
reachable patterns. Choosing the baselines so those patterns are equiprobable gives
$H(A)\to\log_2(R+1)$ bits. Crucially, by \eqref{eq:zbar} each $\Ebase^{(r)}$ sets
$\bar z_r$ directly, so this ladder is obtained without moving any receptor and without
sacrificing any lever arm.

\paragraph{Step 12e. The cost of staggering when $\nu$ is small.} Repeating Steps 9a to
9c with a nonzero mean, $q_r=\Phibar(\mu_r+\nu_r z)$ with $\mu_r=\EE[\tau_r]/\sigma$,
\begin{equation}
  \eta_r=1-\frac{\mu_r^2+\nu_r^2}{\pi\ln2}+O(\cdot^4),
  \label{eq:etar}
\end{equation}
so the first term of \eqref{eq:general} loses $0.459\,\mu_r^2$ bits per receptor while
the second gains at most $\log_2(R+1)$ in total. For large $R$ the loss is extensive and
the gain logarithmic, so staggering is unambiguously bad.

\paragraph{Step 12f. The crossover.} The two arguments meet where the common-mode
standard deviation equals the private-signal standard deviation, which by
\eqref{eq:nudef} is exactly $\nu=1$:
\begin{equation}
  \boxed{\;\text{stagger the thresholds if and only if } \nu>1 .}
\end{equation}
Below the transition the optimal array is a set of identical balanced receptors at the
common radius $U_\star$ of Section~\ref{sec:lever} with spread-out directions. Above it,
the optimal array is a ladder of baselines tiling the range of the common mode. This is
a genuine symmetry-breaking transition in the design, not a smooth reweighting, and it
is the central qualitative prediction of this document.

\begin{remark}[A second staggering axis, opened by a learnable $\Emax$]
Equation~\eqref{eq:sdrad} shows that a receptor's exposure to concentration noise is
set by its dynamic range $\Gamma_r=\Emax^{(r)}/\sigma_c$. Since $\Emax$ is learnable per
unit, the array can also stagger \emph{in $\Gamma$}, dedicating high-$\Gamma$ receptors
to concentration-insensitive identity reading and low-$\Gamma$ receptors to level
reading. This is a prediction with no analogue in a fixed-kernel model and it is easy to
test: inspect the spread of the optimised $\Emax$ values and check whether it correlates
with the receptors' contribution to $I(A;c)$ rather than $I(A;l)$.
\end{remark}

\begin{remark}[Computing the staggered case]
With distinct receptors the pattern probability is no longer a function of $K$, so
\eqref{eq:grouping} does not apply and $I(A;(\|l\|^2,s))$ must be obtained either by
enumerating the $2^R$ patterns at each quadrature node, affordable for $R\lesssim15$, or
with the estimators already implemented in the project. Optimising only the first term
of \eqref{eq:general} is cheap but systematically biased toward the identical design,
since it is blind to precisely the thermometer information that staggering exists to
capture.
\end{remark}

\section{Step 13: the other corner, $R>D$}
\label{sec:corner}

When $R$ exceeds $D$ the directions cannot be orthogonal, Lemma~\ref{lem:orth} gives
$\Cov(p_r,p_{r'})=\sigma^2\uhat_r\cdot\uhat_{r'}\neq0$, and a second source of
redundancy appears that has nothing to do with $\nu$. Conditional independence fails and
the exchangeable reduction is lost. The relevant statement becomes a counting bound: at
fixed concentration each receptive field is a ball \eqref{eq:ball}, $R$ balls in $\RR^D$
lift to $R$ hyperplanes in $\RR^{D+1}$ through $l\mapsto(l,\|l\|^2)$, and the number of
cells of such an arrangement grows polynomially in $R$ with exponent set by the
dimension, so $H(A)=O(D\log_2 R)$. The three corners are therefore
\begin{center}
\begin{tabular}{lll}
\toprule
 & $\nu\lesssim1$ & $\nu\gtrsim1$ \\
\midrule
$R\le D$ & $R$ bits, ideal, flat in every parameter & $\approx R/\nu$ bits, thermometer if staggered \\
$R>D$ & \multicolumn{2}{c}{$O(D\log_2 R)$ bits, geometric saturation} \\
\bottomrule
\end{tabular}
\end{center}
A single figure in the $(R/D,\nu)$ plane with measured MI as colour and these three
asymptotics as annotations would carry the entire analytic section of the paper.

\section{Numerical recipe and validation protocol}
\label{sec:numerics}

\paragraph{The closed form, identical receptors.} Given
$(D,\sigma,\lambda,\Emax,\sigma_c,R)$:
\begin{enumerate}
  \item $\rho^2=D\sigma^2/\lambda^2$, $\Gamma=\Emax/\sigma_c$,
        $P=\rho^2/2$, $Q=De^{2\rho^2}/(4\rho^2\Gamma^2)$.
  \item $\Lambda=W\big(P/(Qe)\big)$, $w_\star=(1+\Lambda)/2$,
        $U_\star=\lambda\sqrt{w_\star}$, $\nu_{\min}=\rho/\sqrt{\Lambda}$.
  \item Check self-consistency with \eqref{eq:selfcons}: if
        $\sqrt{\Lambda D/(2\rho^4)}\lesssim2$, abandon the delta method and use the exact
        route below.
  \item $\Ebase$ follows from $\bar z=e^{-(\rho^2+w_\star)}$ and \eqref{eq:Zlaw}.
  \item On Gauss--Hermite nodes $z_i$ with weights $\omega_i$, set
        $q_i=\Phibar(\nu z_i)$, then
        $\Pi_k=\sum_i\omega_i q_i^k(1-q_i)^{R-k}$ evaluated in log space with
        \texttt{lgamma} for the binomials.
  \item $\PP(K=k)=\binom{R}{k}\Pi_k$. Assert it sums to one, which is the check of
        Step~4d.
  \item $\mathrm{MI}=H(K)+\sum_k\PP(K=k)\log_2\binom{R}{k}$ from \eqref{eq:grouping},
        and independently $R\eta(\nu)+I(A;\tau)$ from \eqref{eq:chain}. The two must
        agree. Disagreement means the quadrature is too coarse in the tails, which is
        where $q$ saturates.
\end{enumerate}

\paragraph{The exact route, without A7.} Do not expand $\ln Z$. Draw or quadrature the
pair $(\|l\|^2,s)$, with $\|l\|^2\sim\sigma^2\chi^2_D$ and $s$ normal, and for each node
compute the demand $Z$ from \eqref{eq:Zlaw} and set
\begin{equation}
  q = \begin{cases}
    1 & Z\le0 \quad\text{(saturated on)}\\
    \Phibar\big(\tau/\sigma\big), \ \ \tau=\dfrac{\|l\|^2+\lambda^2\ln Z+U^2}{2U}
      & 0<Z\le1\\
    0 & Z>1 \quad\text{(silent)} .
  \end{cases}
\end{equation}
This handles the atoms of Step~1d exactly, costs one two-dimensional quadrature, and is
the reference against which the closed form should be checked. Optimise $U$ numerically
on this object rather than trusting \eqref{eq:numin} whenever $\rho\gtrsim1$.

\paragraph{Validating against the simulator.} Run the pipeline with a single family, one
ligand per sniff, \texttt{affinity\_kernel="gaussian"}, a deterministic readout, and
$R\le D$. Sweep $\lambda$ so that $\nu_{\min}$ runs from about $0.3$ to about $3$, which
by Section~\ref{sec:limits} is a sweep of a single scalar. Four predictions are tested
at once, each failure mode diagnostic:
\begin{itemize}
  \item measured MI follows $R\,\eta(\nu_{\min})$, with no free parameter;
  \item optimised receptor norms concentrate near
        $U_\star=\lambda\sqrt{w_\star}$, which is a few times $\lambda$ and \emph{not}
        the cloud radius $\sigma\sqrt D$;
  \item optimised baselines satisfy $\bar z=e^{-(\rho^2+w_\star)}$;
  \item the baselines are tight below $\nu=1$ and spread into a ladder above it.
\end{itemize}
Per the project's standing rule, every point must be averaged over independently drawn
environments, never over evaluation repeats of one model, and the reported spread must
be the spread across independent optimisations.

\section{Summary of results}

\begin{center}
\begin{tabular}{lll}
\toprule
Result & Equation & Status \\
\midrule
Firing condition: match exceeds demand & \eqref{eq:demand} & exact, gaussian kernel \\
Receptive field is exactly a ball & \eqref{eq:ball} & exact, gaussian kernel \\
$\Ebase$ sets the operating point, $\Emax$ the dynamic range & \eqref{eq:Zlaw} & exact \\
Common mode $\Theta=\|l\|^2-\mathrm{Rad}^2(s)$ & \eqref{eq:Theta} & exact \\
Concentration spread $\equiv$ radial cloud breathing & \eqref{eq:Thetastats} & exact \\
Array $=$ $R$ Gaussians cut by one shared threshold & \eqref{eq:tau} & exact under A5, A6 \\
Balance: $\EE[\mathrm{Rad}^2]=D\sigma^2+U^2$ & \eqref{eq:balance} & exact, any kernel \\
Operating point $\bar z=e^{-(\rho^2+w)}$ & \eqref{eq:zbar} & delta method \\
Concentration noise $\propto e^{\rho^2+w}$ & \eqref{eq:sdrad} & delta method \\
Control parameter $\nu(w)$ & \eqref{eq:nugauss} & delta method \\
Quadratic kernel has no optimal $U$ & \eqref{eq:nuquad} & exact in that limit \\
Optimal lever arm, Lambert $W$ & \eqref{eq:numin} & delta method \\
$U_\star\simeq\lambda\sqrt{\ln\Gamma}$, $\nu_{\min}\simeq\rho/\sqrt{2\ln\Gamma}$ & \eqref{eq:scaling} & asymptotic \\
Self-consistency $\Gamma\bar z|_\star=\sqrt{\Lambda D/2\rho^4}$ & \eqref{eq:selfcons} & exact given A7 \\
$\mathrm{MI}=H(K)+\EE[\log_2\binom{R}{K}]$ & \eqref{eq:grouping} & exact under A4 \\
$\eta(\nu)=1-0.459\nu^2$ & \eqref{eq:smallnu} & small $\nu$ \\
$\eta(1)=1/(2\ln2)=0.7213$ & \eqref{eq:nu1} & exact \\
$\eta(\nu)\approx1.040/\nu$ & \eqref{eq:largenu} & large $\nu$ \\
Error of A4 & $O(1/D)$ per receptor & bounded \\
$H(A\mid l)\le\log_2(R+1)$, thermometer & \eqref{eq:split} & exact \\
Stagger if and only if $\nu>1$ & Step~12f & asymptotic argument \\
\bottomrule
\end{tabular}
\end{center}

\paragraph{Open items, in the order they should be attacked.}
\begin{enumerate}
  \item Implement $\eta(\nu)$ and the exact MI, verify the three limits and the exact
    point $\eta(1)=1/(2\ln2)$.
  \item Implement the exact route of Section~\ref{sec:numerics} and compare with the
    Lambert-$W$ optimum over a grid in $\rho$, to map where A7 is trustworthy.
  \item Validate against the simulator along the $\lambda$ sweep, checking the four
    predictions listed above.
  \item Optimise \eqref{eq:general} over baselines and norms, and locate the staggering
    transition numerically.
  \item Only then, the correlated-direction case $R>D$, and the heteromer questions of
    Appendix~\ref{app:hetero}.
\end{enumerate}

\appendix

\section{Heteromers}
\label{app:hetero}

Nothing in the main text depends on this appendix. It is recorded here because the
results are short, exact, and decide what the heteromer comparison is allowed to claim.
Throughout, a heteromer is assembled from $\ksub$ subunits with latent coordinates
$v_1,\dots,v_{\ksub}$, baselines $\Ebase^{(u)}$ and ceilings $\Emax^{(u)}$, and its
half-activation energy is the arithmetic mean of the subunit energies.

\subsection{The match of a heteromer is the weighted mean of its subunit matches}

Average the gaussian kernel \eqref{eq:kernel} over the subunits and collect terms:
\begin{align}
  \ln EC_{50}^{(\mathrm{het},l)}
   &= \frac{1}{\ksub}\sum_{u}\left[\Ebase^{(u)}
      +\Emax^{(u)}\Big(1-e^{-\|l-v_u\|^2/\lambda^2}\Big)\right] \notag\\
   &= \bar{E}_{\mathrm{base}} + \bar{E}_{\mathrm{max}}
      - \frac{1}{\ksub}\sum_u \Emax^{(u)} e^{-\|l-v_u\|^2/\lambda^2},
\end{align}
with $\bar{E}_{\mathrm{base}}$ and $\bar{E}_{\mathrm{max}}$ the subunit means. Defining
the \textbf{pooled match}
\begin{equation}
  K_{\mathrm{het}} \;\equiv\;
   \frac{1}{\ksub\,\bar{E}_{\mathrm{max}}}\sum_u \Emax^{(u)}\,
   e^{-\|l-v_u\|^2/\lambda^2}
   \;\in\;(0,1],
\end{equation}
the whole expression takes the homomer form,
$\ln EC_{50}^{(\mathrm{het})}=\bar{E}_{\mathrm{base}}
 +\bar{E}_{\mathrm{max}}(1-K_{\mathrm{het}})$, so by Step~1b
\begin{equation}
  \boxed{\;a_{\mathrm{het}}=\mathbf{1}\big[K_{\mathrm{het}}>Z_{\mathrm{het}}\big],
  \qquad
  Z_{\mathrm{het}}=\frac{\bar{E}_{\mathrm{base}}+\bar{E}_{\mathrm{max}}-s}
                        {\bar{E}_{\mathrm{max}}} . }
\end{equation}
The match-versus-demand picture therefore survives heteromerisation intact. The only
change, and it is the important one, is that the match is now a weighted sum of
$\ksub$ Gaussian bumps instead of one.

\subsection{Consequence: the field is not a ball}

Because a sum of Gaussians is not a Gaussian, the firing set
$\{l : K_{\mathrm{het}}(l)>Z\}$ is a superlevel set of a Gaussian mixture, not a ball.
A single perfectly matched subunit contributes only $\Emax^{(u)}/(\ksub\bar{E}_{\max})
\approx 1/\ksub$ to the pooled match, so the logic of the receptor depends on where its
demand sits:
\begin{itemize}
  \item $\bar z\lesssim1/\ksub$: one well-matched subunit is enough, the heteromer acts
    as a soft OR over its subunits and its field is a union of blobs;
  \item $\bar z\gtrsim1/\ksub$: several subunits must engage at once, the heteromer acts
    as a soft AND and its field is concentrated where the bumps overlap.
\end{itemize}
With the balanced operating point $\bar z=e^{-(\rho^2+w)}$ of \eqref{eq:zbar} and the
worked numbers of Step~7e, $\bar z\approx0.34$ against $1/\ksub=0.2$, so the model sits
on the AND side but not far from the boundary. \textbf{This is the only mechanism by
which a heteromer can do something no homomer can do at fixed receptor count}, and
A.4 below shows that it is the only one that survives scrutiny.

\subsection{The harmonic limit collapses the field back to a ball}

Expand $e^{-d^2/\lambda^2}\simeq1-d^2/\lambda^2$ with common $\Emax$, and apply the
parallel-axis identity with $\bar v=\frac{1}{\ksub}\sum_u v_u$ and
$S^2=\frac{1}{\ksub}\sum_u\|v_u-\bar v\|^2$:
\begin{equation}
  \frac{1}{\ksub}\sum_u\|l-v_u\|^2=\|l-\bar v\|^2+S^2 ,
\end{equation}
which follows by writing $l-v_u=(l-\bar v)-(v_u-\bar v)$, squaring, and noting that the
cross term vanishes because $\sum_u(v_u-\bar v)=0$. Hence
\begin{equation}
  \boxed{\;\ln EC_{50}^{(\mathrm{het},l)}
   = \bar{E}_{\mathrm{base}} + k\big(\|l-\bar v\|^2+S^2\big),
   \qquad k=\frac{\Emax}{\lambda^2}. }
  \label{eq:hetharm}
\end{equation}
\textbf{In the harmonic limit a heteromer is exactly a homomer placed at the centroid of
its subunits, with its baseline raised by the subunit dispersion.} Heteromerisation is
then not a richer functional form at all, only a reparametrisation.

\subsection{The centroid theorem, and why it forbids a fixed-$R$ advantage}

Give each gene the pair $P_u=(v_u,\,B_u)$ with
$B_u\equiv k\|v_u\|^2+\Ebase^{(u)}$ its offset in the sense of \eqref{eq:Theta}. The
offset of a heteromer follows from \eqref{eq:hetharm} and the parallel-axis identity
once more:
\begin{equation}
  B_{\mathrm{het}} = k\big(\|\bar v\|^2+S^2\big)+\bar{E}_{\mathrm{base}}
   = \frac{1}{\ksub}\sum_u\Big[k\|v_u\|^2+\Ebase^{(u)}\Big]
   = \frac{1}{\ksub}\sum_u B_u ,
\end{equation}
and of course $v_{\mathrm{het}}=\bar v$. Therefore
\begin{equation}
  \boxed{\;\big(v,B\big)_{\mathrm{het}}
   = \text{centroid of } \big\{(v,B)_u : u \in \text{subunits}\big\}. }
\end{equation}
This statement is true \emph{only because $\Ebase$ is learnable}: it is what makes $B_u$
a free coordinate rather than a function of $v_u$. Two corollaries follow.

\paragraph{Containment.} A homomer of gene $u$ realises an arbitrary point
$(v_u,B_u)\in\RR^{D+1}$, since position and baseline are independently learnable. A
heteromer realises a centroid of such points, which is also a point of $\RR^{D+1}$ and
nothing more. Hence the set of $R$-tuples achievable by an array of $R$ homomers
\emph{contains} the set achievable by an array of $R$ heteromers, and
\begin{equation}
  \boxed{\;\mathrm{MI}^{*}_{\mathrm{homomer}}(R)\;\ge\;\mathrm{MI}^{*}_{\mathrm{heteromer}}(R)
  \qquad\text{exactly, in the harmonic limit.}}
\end{equation}
\textbf{Comparing the two architectures at fixed receptor count is therefore
meaningless: homomers win by construction.} Any heteromer advantage must come from
either a fixed \emph{gene} budget, or the non-ball fields of A.2.

\paragraph{The honest comparison.} With $n$ genes a homomer array has $n$ receptors,
while heteromers give $\binom{n+\ksub-1}{\ksub}$ of them. For the octopus $n=26$ and
$\ksub=5$, so the heteromer repertoire is of order $10^5$ against $26$. The question the
paper can legitimately ask is whether a large, contracted, correlated set beats a small,
free one, and the answer will be governed by whether the array is limited by receptor
count, in which case heteromers win, or by the common mode and the geometry of
Section~\ref{sec:corner}, in which case they do not.

\subsection{The lever-arm penalty at fixed gene positions}

With free baselines, balance fixes $B$ and leaves the lever arm as the only difference
between architectures. By the triangle inequality $\|\bar v\|\le\frac{1}{\ksub}\sum_u\|v_u\|$,
with equality only if all subunits are parallel, and for subunits of common norm $G$ in
random directions $\|\bar v\|\approx G/\sqrt{\ksub}$. Since
$\nu\propto1/U$ by \eqref{eq:nudef},
\begin{equation}
  \frac{\nu_{\mathrm{het}}}{\nu_{\mathrm{hom}}}\;\approx\;\sqrt{\ksub}\;\approx\;2.2
  \qquad(\ksub=5).
\end{equation}
There is also a design conflict worth stating: by Section~\ref{sec:lever} the optimal
lever arm is $U_\star$, but one gene cloud cannot give $U_\star$ to both its homomers and
its heteromers, since the latter sit a factor $\sqrt{\ksub}$ closer in. A gene cloud
optimised for its heteromers places its homomers too far out, and conversely. Quantifying
that compromise, which requires the joint distribution of $(\bar v,S^2)$ over subunit
multisets, is the natural next calculation and is not attempted here.


\end{document}
